\scanpage{222}
1) Soit $f:X\longrightarrow S$ morphisme lisse d'espaces analytiques (lisses)

$\omega$ une 1-forme fermée \inserted{(holomorphe)} relative sur $X/S$

$a$ et $b$ deux sections \inserted{holomorphes} de $X$ sur $S$, telles que $\forall s\in S$, $a(s)$ et $b(s)$ \inserted{sont dans la \unsure{$=$|même}. comp. connexe de $X_s$}.

Localement sur $S$, on peut trouver un arc \deleted{diff.} \inserted{différentiable par morceaux} $C_s:[0,1]=I\longrightarrow X_s$, joignant $a(s)$ et $b(s)$, et dépendant continûment de $s$.

Utilisant le fait que $\omega$ est fermée, donc localement de la forme $d_{X/S}\varphi$, on constate que la fonction sur $S$
\[
I(s)=\int_{C_s}^{\deleted{\unsure{b(s)}}}\omega_s=\int_{a(s)}^{b(s)}\omega_s
\]
est holomorphe en $s$. [On \inserted{\unsure{se place} au voisinage de $s_0\in S$ et \unsure{se}} ramène au cas où \deleted{\unsure{a(s_0)\text{ et}}} $a(s_0)$ et $b(s_0)$ sont assez voisins pour que dans un ouvert $U$ contenant $a(s_0),b(s_0)$, on ait $\omega|U=d\varphi$, donc pour $s$ voisin de $s_0$, $I(s)=\varphi b(s)-\varphi a(s)$].

Changeant le système des $C_s$ revient à le remplacer par $C'_s=C_s+Z_s$, $Z_s$ un 1-cycle, donc on trouve
\[
(1)\qquad I'(s)=I(s)+\underbrace{\int_{Z_s}\omega_s}_{\deleted{\unsure{\Pi}}}
\]

Considérons une expression formelle (avec des $\omega_i$ comme $\omega$, et des op. diff. \inserted{(hol.)} $D_i$ sur $S$]
\[
(2)\qquad \sum D_{is}\int_{a(s)}^{b(s)}\omega_{is}
\]
\margin{221}
